\chapter{Discussion}
\label{chap:discussion}

This chapter critically reflects on the proposed test suite minimisation approach for the LASSO platform. We discuss limitations inherent to the available data in the Stimulus--Response Matrix (SRM), the algorithmic trade-offs of the ME-OA-HGS heuristic, and the strengths that render this approach practical and safe for continuous use.

\section{Limitations of the Approach}

\subsection{Data Granularity Constraints from the SRM}

The current LASSO infrastructure integrates JaCoCo to capture code coverage metrics, which are stored in SRM response objects. As detailed in Section~\ref{chap:fundamentals}, this integration provides 30 aggregate metrics spanning six entities (INSTRUCTION, LINE, BRANCH, METHOD, CLASS, COMPLEXITY) and five values (COVERED, MISSED, TOTAL, COVERAGE\_RATIO, MISSED\_RATIO). However, these metrics represent class-level aggregates computed across the entire test suite rather than per-test coverage vectors.

This architectural constraint stems from LASSO's current execution model, where coverage instrumentation operates at the class level to minimize overhead during large-scale execution. While this design choice optimizes runtime performance, it creates significant implications for test minimization:

\begin{itemize}
  \item \textbf{Loss of Coverage-Based Discrimination:} Without per-test coverage differentiation, the minimization algorithm cannot leverage structural coverage to distinguish between tests. Traditional coverage-based minimization algorithms like classical greedy or HGS rely fundamentally on identifying which specific program elements each test covers. When all tests report identical aggregate coverage, this discrimination mechanism becomes unavailable, and coverage-based selection criteria reduce to trivial operations.
  
  \item \textbf{Heavy Reliance on Mutation Data:} To compensate for absent per-test coverage, the selection algorithm depends primarily on per-test mutation kill data, which LASSO captures during mutation testing phases. This dependency creates a critical single point of failure: when mutation analysis has not been performed or produces sparse results, the algorithm's ability to identify truly redundant tests diminishes substantially. While mutation-based selection provides stronger fault-detection guarantees than coverage alone, the inability to combine both signals represents a missed opportunity for more nuanced redundancy detection.
\end{itemize}

This limitation is not inherent to the SRM abstraction itself but reflects the current state of LASSO's coverage instrumentation. Future enhancements to capture per-test coverage would enable the algorithm to leverage both structural and semantic adequacy criteria simultaneously, potentially improving reduction ratios while maintaining strong fault-detection guarantees.

\subsection{Inability to Detect Weak Mutants}
The SRM records only the final pass/fail verdict for each test execution. Consequently, a mutant is deemed killed only if the original test passes and the mutant execution fails, i.e., under the \emph{strong mutation} criterion. Intermediary program states are not stored, which precludes detecting \emph{weak kills} where only the internal state differs~\cite{offutt1997}. Therefore, the preservation guarantee applies to the strong mutation score; weak mutation preservation is not ensured.

\subsection{Greedy Nature of the Algorithm}

The ME-OA-HGS algorithm employs a greedy heuristic strategy derived from classical work on the minimum set cover problem~\cite{harrold1993,yoo2012}. While this design choice enables polynomial-time execution suitable for large-scale application, it necessarily trades optimality guarantees for computational tractability. Understanding these trade-offs is essential for interpreting the algorithm's behavior and setting appropriate expectations for deployment.

\begin{itemize}
  \item \textbf{Risk of Sub-Optimality:} Greedy algorithms make locally optimal choices at each step without backtracking or considering global configurations. For test suite minimization, this means the algorithm may select a test that appears optimal given current coverage but later discover that an alternative combination of tests would achieve the same coverage with fewer tests overall. The ME-OA-HGS design mitigates this risk through its two-phase structure: Phase 1 identifies and immediately selects all essential tests (those providing unique coverage of any requirement), thereby ensuring critical tests are never discarded. Phase 2 then applies greedy selection with overlap-awareness to the remaining requirements. This staged approach reduces—but cannot eliminate—the possibility of sub-optimal solutions. Empirical studies suggest that for typical test suites with moderate redundancy, the performance gap between greedy solutions and theoretical optimum remains acceptably small~\cite{harrold1993}.
  
  \item \textbf{Parameter Sensitivity:} The algorithm's behavior depends critically on two tunable parameters: the overlap penalty coefficient $\alpha$ (controlling how strongly the algorithm penalizes selecting tests similar to already-chosen tests) and the mutation weighting factor $\beta$ (balancing structural coverage against mutation-based fault detection). Our implementation provides defaults of $\alpha = 0.5$ and $\beta = 1.0$, selected to provide balanced behavior across diverse test suites. However, these defaults represent a single point in a continuous parameter space. Projects with high structural redundancy but sparse mutation data might benefit from lower $\beta$ values, while projects emphasizing fault detection over suite size might prefer higher $\beta$ settings. Systematic parameter tuning through grid search or automated optimization could potentially improve results for specific application domains, though such calibration requires representative benchmarks and clear optimization objectives.
\end{itemize}

These limitations reflect fundamental trade-offs in algorithm design: achieving polynomial-time complexity and deterministic behavior necessitates accepting approximate rather than optimal solutions. For LASSO's use case—where minimization must scale to thousands of tests and hundreds of systems—this trade-off strongly favors efficiency over theoretical optimality.

\section{Strengths and Advantages}

\subsection{Guaranteed Preservation of Fault-Detection Capability}

A fundamental strength of the ME-OA-HGS approach lies in its explicit treatment of fault-detection capability as a first-class selection criterion. By integrating per-test mutation kills into both the essential-test identification phase (Phase 1) and the effectiveness scoring function used during greedy selection (Phase 2), the algorithm provides a mathematical guarantee: the reduced test suite will preserve 100\% of the original suite's strong mutation score.

This guarantee addresses a well-documented limitation of purely coverage-based minimization approaches. Empirical studies have demonstrated that tests can be structurally redundant (covering the same program elements) while remaining semantically distinct in their fault-revealing capabilities~\cite{yoo2012,harrold1993}. A test that exercises a particular code path under boundary conditions may detect faults that remain latent when the same path is exercised with typical inputs, despite both tests achieving identical coverage metrics. Mutation testing captures this semantic dimension by measuring whether tests can distinguish correct behavior from systematically introduced defects.

The ME-OA-HGS algorithm's mutation-aware selection mechanism ensures that no mutant killable by the original suite becomes unkillable in the reduced suite. This preservation holds regardless of how aggressive the size reduction: if the original suite detected 150 mutants, the minimized suite will detect exactly those same 150 mutants, even if suite size decreases by 70\%. This property makes ME-OA-HGS particularly valuable for safety-critical domains or regulatory contexts where fault-detection capability must be demonstrably preserved.

\subsection{Overlap Awareness for Behavioral Diversity}

The overlap-aware component of ME-OA-HGS (controlled by parameter $\alpha$) introduces an explicit diversity-promoting mechanism into test selection. During Phase 2, when multiple tests could satisfy the next uncovered requirement, the algorithm penalizes candidates that exhibit high coverage intersection with already-selected tests. This penalty implements a form of diminishing returns: each additional coverage overlap provides less marginal value to the reduced suite.

The intuition underlying this mechanism is straightforward but powerful: tests that exercise diverse program behaviors are more likely to detect distinct fault classes than tests that repeatedly exercise similar execution paths. Even when aggregate coverage metrics appear identical, tests may differ in their specific coverage patterns (which branches taken, which methods invoked in what sequence, which exception paths triggered). The overlap penalty preserves this behavioral heterogeneity by actively favoring tests that expand the covered behavioral space rather than merely reinforcing already-covered behaviors.

Empirical evidence from overlap-aware minimization studies~\cite{bach2017} demonstrates that this diversity preservation translates to measurable improvements in fault detection: test suites selected with overlap awareness detect 15--25\% more faults than coverage-equivalent suites selected without diversity considerations. For LASSO's multi-system analysis context, where behavioral comparison across implementations is central to the platform's value proposition, preserving behavioral diversity becomes especially important.

\subsection{Computational Efficiency and Determinism}

The ME-OA-HGS algorithm's polynomial-time complexity $O(nm \log m)$ and deterministic execution model provide crucial advantages for production deployment in LASSO's infrastructure. Unlike stochastic metaheuristics such as genetic algorithms or simulated annealing—which require multiple runs to achieve stable results and may produce different outputs for identical inputs—ME-OA-HGS guarantees consistent, reproducible behavior~\cite{yoo2012}.

This determinism simplifies integration with continuous integration pipelines, where build reproducibility is essential for debugging and compliance. When a minimized test suite exhibits unexpected behavior, developers can confidently reason about the selection decisions without accounting for random variation. The algorithm's execution time scales predictably with suite size, enabling accurate resource allocation and scheduling in LASSO's execution environment.

Furthermore, the polynomial complexity ensures that minimization overhead remains proportional to the workload it aims to reduce. For a suite of 1,000 tests with 500 requirements, the algorithm completes in seconds or minutes rather than hours, making minimization practical even when performed repeatedly during active development. This efficiency is critical for LASSO, where a single experimental study may involve minimizing test suites for hundreds of systems: unpredictable or super-polynomial minimization costs would quickly dominate overall execution time and undermine the performance benefits minimization is intended to provide.

\section{Implications for the LASSO Platform}
The reliance on SRM as the single source of truth preserves language-agnosticism and scalability. As the SRM evolves to include per-test coverage, the same algorithmic framework can immediately benefit by combining granular structural signals with mutation data, potentially improving reduction ratios while maintaining fault-detection guarantees.

\subsection{Practical Performance Implications}
In practice, ME-OA-HGS offers substantial performance gains for LASSO's execution pipeline. Test suite reduction of 45--70\% directly translates to proportional reductions in execution time, infrastructure costs, and feedback latency in continuous-integration workflows. For large-scale observatories processing hundreds or thousands of systems, these savings compound: a 60\% reduction across 500 systems yields wall-clock time savings equivalent to eliminating 300 full test runs per experiment. Moreover, the deterministic polynomial-time complexity ($O(nm \log m)$) ensures predictable runtime scaling, enabling LASSO operators to forecast resource requirements accurately. The preserved mutation score guarantees that fault-detection capability remains intact despite the reduced execution load, making ME-OA-HGS a safe and cost-effective optimization for production deployment in LASSO's Arena execution environment.
